\section{Embedded root systems and moment-certified sections}\label{sec:sections}
An embedded root system imposes restrictions on the projections of the
entire minimal shell. In the case below, these restrictions determine
the number of points perpendicular to a child section.

\begin{theorem}[$E_8$-extension count]\label{d43:thm:structure}
Let $L$ be an extremal even unimodular lattice of rank 48, and let
$\iota:\sqrt3E_8\hookrightarrow L$ be an isometric embedding.
Use the simple-root numbering in~\eqref{d43:eq:simple}, and set
\[
 U=\operatorname{span}_{\R}\{\iota(\sqrt3\alpha_1),\ldots,
                              \iota(\sqrt3\alpha_5)\}.
\]
Then exactly $2,553,792$ vectors of squared norm 6 in $L$ are
perpendicular to $U$. The same count holds for every $W(E_8)$-conjugate
of this $D_5$ subsystem inside the embedded $E_8$.
\end{theorem}

The proof uses integrality, the lattice minimum and spherical-design
moments, not an extension of the Weyl group to automorphisms of $L$.

\subsection{Signature domains and target statistics}
The standard design identities~\cite{dgs,nebe-designs} give the following
principle. The signature domain must be complete, and averaging must
preserve the target statistic.

\begin{proposition}[Finite-section statistics]\label{prop:section-statistics}
Let $L\subset\R^n$ be an integral lattice of minimum squared norm $\mu$,
and suppose its minimal shell $X$ has $N$ points and is a spherical
$(2s+1)$-design. Choose independent $b_1,\ldots,b_r\in L$ with Gram
matrix $H$. For $x\in X$ set
\[
y(x)=(\ip{x}{b_i})_{i=1}^r,\qquad
q(y)=y^tH^{-1}y,\qquad f(y)=|\{x\in X:y(x)=y\}|.
\]
Choose a finite antipodal set $\mathcal R\subset\Z^r\setminus\{0\}$
of section-vector coefficients, and put
\begin{equation}\label{eq:section-domain}
\begin{split}
E&=\{Hz:z\in\mathcal R,\ z^tHz=\mu\},\\
D&=\{y\in\Z^r:q(y)\le\mu,\quad
             |z^ty|\le z^tHz/2\ \text{for every }z\in\mathcal R\}.
\end{split}
\end{equation}
Then every signature lies in the disjoint union $D\cup E$, and
$f(e)=1$ for each $e\in E$. No primitivity of the section is required.

For a multi-index $\alpha$ of total degree at most $2s+1$, define
$M_\alpha=0$ when $|\alpha|$ is odd, and, when $|\alpha|=2m$, define
\begin{equation}\label{eq:section-moments}
M_\alpha=
\frac{N\mu^m}{\prod_{j=0}^{m-1}(n+2j)}
\sum_P\prod_{(a,b)\in P}H_{i_a i_b},
\end{equation}
where the labelled list $(i_1,\ldots,i_{2m})$ contains $\alpha_i$
copies of $i$, and $P$ runs over its pairings. Empty products give
$M_0=N$.

Let $\Gamma$ be a finite group of linear transformations preserving
$q$, $D$ and $E$, and let $\psi:D\cup E\to\R$ be $\Gamma$-invariant.
For its orbits $O_j$ on $D$, form
\[
A_{\alpha j}=\sum_{y\in O_j}y^\alpha,\qquad
b_\alpha=M_\alpha-\sum_{e\in E}e^\alpha,\qquad
c_j=\sum_{y\in O_j}\psi(y).
\]
Any selected moment rows in the stated degree range may be used. If
$A^t\lambda\le c$ entrywise, then
\begin{equation}\label{eq:section-statistic-bound}
\sum_{x\in X}\psi(y(x))\ \ge\
\sum_{e\in E}\psi(e)+\lambda^tb.
\end{equation}
If $A^t\lambda=c$, equality holds.
\end{proposition}
\begin{proof}
Integrality and orthogonal projection give $y(x)\in\Z^r$ and
$q(y(x))\le\mu$. For $a=\sum_i z_i b_i$ and $x\ne\pm a$,
$\|x\pm a\|^2\ge\mu$ gives $|z^ty(x)|\le\|a\|^2/2$.
The exceptions have $\|a\|^2=\mu$: their full norm lies in the
section, so each fibre is a singleton and violates its own domain
inequality. Thus $D\cap E=\varnothing$; positive definiteness of $q$
makes $D$ finite.

Polarization gives~\eqref{eq:section-moments}. Since every isometry of
$q$ extends orthogonally to $\R^n$, these moments are $\Gamma$-invariant
even when $\Gamma$ does not act on $L$. Hence
\[
\bar f(y)=\frac1{|\Gamma|}\sum_{g\in\Gamma}f(g^{-1}y).
\]
is nonnegative, has the same moments and equals one on $E$.
Its orbit values satisfy $Av=b$, while invariance of $\psi$ preserves
$\sum_X\psi=\sum_E\psi+c^tv$. The result follows from
$(c-A^t\lambda)^tv\ge0$.
\end{proof}

Antipodality removes odd total-degree rows, not monomials with odd
individual exponents. Exceptional orbits may instead be retained with
value one. A noninvariant target is only averaged, not determined for a
fixed child. In dimension 45, longer section vectors remove all columns
carrying slack in a one-sided certificate.

\subsection{A realization in \texorpdfstring{$P_{48n}$}{P48n}}
Take $L=P_{48n}$, whose Gram matrix $G$ is published as CQ48a
in the Nebe--Sloane catalogue~\cite{latticecatalogue}. It is an extremal even
unimodular lattice of rank 48, minimum squared norm 6, with minimal shell
\[
X=\{x\in L:\|x\|^2=6\},\qquad |X|=52,416,000.
\]
The parent certificate~\cite{crosssections} specifies eight integer
coefficient rows $T$ and scaled coordinate rows $B$,
satisfying
\begin{equation}\label{d43:eq:embedding}
 TGT^t=H=3C(E_8)=BB^t/96.
\end{equation}
The catalogue and certificate Grams agree entry by entry. The first five
rows $T_5$ have Gram
\begin{equation}\label{d43:eq:d5}
H_5=\begin{pmatrix}
6&0&-3&0&0\\0&6&0&-3&0\\-3&0&6&-3&0\\
0&-3&-3&6&-3\\0&0&0&-3&6
\end{pmatrix},\qquad \det H_5=972.
\end{equation}
This is $3C(D_5)$ in the stated numbering. Both embeddings are primitive
(all nonzero Smith invariants are one), though the theorem does not
require this. For $U=\operatorname{span}_{\R}(T_5)$, $\dim U^\perp=43$.

\begin{corollary}\label{d43:thm:main}
The set $Y=X\cap U^\perp$ has $2,553,792$ elements. Consequently the ordinary
Euclidean kissing number satisfies $K(43)\ge2,553,792$.
\end{corollary}

Theorem~\ref{d43:thm:structure} gives the cardinality. Since
$\|x-x'\|^2\ge6$ for distinct shell vectors, $Y/\sqrt6$ is a
kissing configuration.

The construction is specified without a list of millions of points. In
integer basis coordinates its defining set is
\begin{equation}\label{d43:eq:integer}
\mathcal Z=\{z\in\Z^{48}:zGz^t=6,\ T_5Gz^t=0\}.
\end{equation}
For $EE^t=G$, its points are $zE/\sqrt6$ with $z\in\mathcal Z$;
the certificate supplies $G,T_5$.

\subsection{The complete projection domain}
Now let $L$ be any lattice satisfying Theorem~\ref{d43:thm:structure},
and let $X$ be its squared-norm-six shell. The theta-series
identity~\eqref{eq:theta48} gives $|X|=52,416,000$.
Identify the parent section with $\sqrt3E_8$, with simple roots
$\sqrt3\alpha_1,\ldots,\sqrt3\alpha_8$. Write the orthogonal projection of
$x\in X$ onto this section as $p=\lambda/\sqrt3$.
Integrality of $L$ and self-duality of $E_8$ imply $\lambda\in E_8$. Set
\[
y_i=\ip{\lambda}{\alpha_i}=\ip{x}{\sqrt3\alpha_i},\qquad
q=\|p\|^2=\|\lambda\|^2/3=y^tH^{-1}y\le6.
\]
For an $E_8$ root $r$, if $x\ne\pm\sqrt3r$, applying the minimum bound to
$x\pm\sqrt3r$ gives
\begin{equation}\label{d43:eq:rootbound}
|\ip{\lambda}{r}|\le3.
\end{equation}
The exceptions are $\lambda=3r$. Each such projection has multiplicity one:
its squared length is already 6, so the perpendicular component vanishes.

For a transparent finite enumeration use the model
\[
E_8=\{u/2:u\in\Z^8,\ u_i\text{ all of the same parity},\
\textstyle\sum_i u_i\equiv0\pmod4\}.
\]
Its roots are the 112 vectors $\pm e_i\pm e_j$ and the 128 half-integer
vectors with entries $\pm1/2$ and an even number of minus signs.
The simple roots used here are
\begin{equation}\label{d43:eq:simple}
\begin{aligned}
\alpha_1&=(1,-1,-1,-1,-1,-1,-1,1)/2, &\alpha_2&=e_1+e_2,\\
\alpha_{i+2}&=-e_i+e_{i+1} &&(1\le i\le6).
\end{aligned}
\end{equation}
Their Gram is precisely $H/3$.

Enumeration of $\|u\|^2\le72$ subject to~\eqref{d43:eq:rootbound} gives
\begin{center}
\begin{tabular}{rrrrrr}\toprule
$\|\lambda\|^2$ &0&2&4&6&8\\\midrule
Nonexceptional signatures &1&240&2,160&6,720&17,280\\\bottomrule
\end{tabular}
\end{center}
There are no other nonexceptional signatures. Including the 240 exceptions
gives 26,641 signatures in total.

For completeness, pair roots bound the two largest absolute coordinates
of $u$ by a sum of 6. Half-integer roots bound $\|u\|_1$ by 12,
except when all coordinates are nonzero with odd sign parity, when they
bound $\|u\|_1-2\min_i|u_i|$ by 12. Thus $\|u\|_1\le18$ is a
safe enumeration bound. Exact parity, norm and root tests give the table;
at norm eight they remove the 240 doubled roots from 17,520 points.

\subsection{Moments and the averaging explanation}
Venkov's theorem~\cite[Theorem 4.1]{nebe-designs} makes $X$ an 11-design,
so~\eqref{eq:section-moments} applies through degree ten. The abstract
Weyl group has five nonexceptional orbits, with sizes $s_j$ and projected
squared norms $q_j=0,2/3,4/3,2,8/3$. Their averaged multiplicities $a_j$
satisfy the radial equations
\begin{equation}\label{d43:eq:radial}
\sum_j s_ja_jq_j^k=|X|6^k\frac{(4)_k}{(24)_k}-240\,6^k,
\qquad 0\le k\le5,
\end{equation}
where $(a)_k$ is the rising factorial and the subtraction removes the
240 section roots. The first five equations give
\[
(a_0,a_1,a_2,a_3,a_4)
=(1,092,000,116,235,9,180,495,63/4).
\]
The equation for $k=5$ also holds.

The roots perpendicular to the first five simple roots form $A_3$.
The target subspace contains respectively $1,12,6,24,0$ admissible
signatures at the five norms, and 12 exceptional roots. Thus the average
contact count over the Weyl images of this $D_5$ child is
\begin{equation}\label{d43:eq:count}
1,092,000+12(116,235)+6(9,180)+24(495)+12=2,553,792.
\end{equation}
This proves an average count, not yet the count of a fixed child.
The $a_j$ need not be integers because they are orbit averages.

\subsection{Exact count for the specified child}
For the fixed child use the target-preserving subgroup
\[
\Gamma=\langle W(D_5),W(A_3),-1\rangle.
\]
This group preserves the full domain and the condition
$y_1=\cdots=y_5=0$, and has 43 orbits $O_j$.
For averaged multiplicities $v_j$, use rows
\[
A_{\alpha,j}=\sum_{y\in O_j}y^\alpha,
\]
with right-hand side~\eqref{eq:section-moments}, and fix exceptional
orbit values to one. The target is
\[
c_j=|\{y\in O_j:y_1=\cdots=y_5=0\}|.
\]
The supplied certificate selects 43 moment or fixed-root rows, and provides
a rational column vector $w$ such that
\begin{equation}\label{d43:eq:dual}
A^tw=c,\qquad w^tb=2,553,792.
\end{equation}
Thus $|X\cap U^\perp|=c^tv=w^tb$. All inputs depend only on the
root system and spherical moments. Precomposing $\iota$ with a Weyl
element proves the conjugate cases without extending it to $L$.
The verifier regenerates $A,b,c$ from the coordinate model, taking only
the row selection and $w$ from the certificate in
Table~\ref{tab:certificates}.

\subsection{A nested family of fixed coordinate sections}
Extend $\iota$ linearly and define the eight fixed coordinate cuts
\[
 U_k=\iota\bigl(\sqrt3\operatorname{span}\{e_1,\ldots,e_k\}\bigr),
 \qquad 1\le k\le8.
\]
For $k\ge2$ these are the coordinate $D_k$ spans ($D_2=A_1^2$,
$D_3=A_3$); $U_1$ is an axis, not a $D_1$ root subsystem, and
$U_8$ is the whole parent.

\begin{theorem}[Coordinate-section family]\label{thm:e8-coordinate-sections}
Let $L,X,\iota$ satisfy Theorem~\ref{d43:thm:structure}.
For $m=8-k$, the number of minimal vectors perpendicular to the fixed
space $U_k$ is
\begin{equation}\label{eq:e8-coordinate-count}
\begin{split}
 |X\cap U_k^\perp|={}&1092000+18360m+464944\binom m2
                  +11880\binom m3\\
                 &+146880\binom m4+2520\binom m5+31680\binom m6.
\end{split}
\end{equation}
The same holds for every Weyl conjugate of this nested family.
In particular, a Weyl conjugate contains the specified $D_5$ child
of Theorem~\ref{d43:thm:structure} as its rank-five member.
\end{theorem}

\begin{proof}
Use the doubled coordinates $u=2\lambda$ of the complete
nonexceptional domain $D$ above. Let $\Gamma_k$ consist of coordinate
permutations within the first $k$ and last $8-k$ positions and sign
changes in an even number of positions. This subgroup of $W(D_8)$
preserves $D$ and the target $u_1=\cdots=u_k=0$. Its orbits are
specified by the two multisets of squared coordinates and $\prod_i u_i$.
The product distinguishes sign parity when no coordinate is zero;
a zero coordinate removes that restriction.

For a partition $a$, write $m_a$ for the monomial symmetric polynomial
whose distinct monomials have coefficient one. Use all rows
\[
 F_{a,b}(u)=m_a(u_1^2,\ldots,u_k^2)
           m_b(u_{k+1}^2,\ldots,u_8^2),\qquad
 |a|+|b|\le5,
\]
where the partition lengths are at most $k$ and $8-k$, respectively.
For orbit masses $w_j$ set $A_{(a,b),j}=F_{a,b}(u_j)$, with
$u_j\in O_j$. Subtracting the 240 singleton exceptions gives $Aw=b$
for the actual masses, without assuming fibre symmetry. The functionals
$u_i$ have Gram $12I_8$, so for $|a|+|b|=d$ the unadjusted moment is
\begin{equation}\label{eq:e8-coordinate-moment}
 \frac{52416000\,72^d}{48\cdot50\cdots(48+2d-2)}
 \nu_k(a)\nu_{8-k}(b)
 \prod_i(2a_i-1)!!\prod_j(2b_j-1)!!,
\end{equation}
where $\nu_l(a)$ is the number of distinct permutations of $a$
padded with zeros to length $l$.

All eight matrices have full column rank: regenerated square minors
have the nonzero residues in Table~\ref{tab:e8-coordinate-family}.
Thus $w$ is unique over $\Q$. The candidate
\[
 w_j=|O_j|a_r,\qquad
 (a_0,a_2,a_4,a_6,a_8)=(1092000,116235,9180,495,63/4),
 \quad r=\|u_j\|^2/4,
\]
comes from the full Weyl average and satisfies every row exactly.
It must therefore equal the actual orbit masses. This does not assign
the fractional value $63/4$ to individual fibres.

It remains to count the target signatures. Since $k\ge1$, a vector
$\lambda$ perpendicular to $U_k$ has integer coordinates and even
coordinate sum in the remaining $m\le7$ positions. At norms
$0,2,4,6,8$, the numbers of admissible signatures are respectively
\[
 1,\quad 4\binom m2,\quad 2m+16\binom m4,\quad
 24\binom m3+64\binom m6,\quad 160\binom m5.
\]
These come from the signed shapes
$0$, $(1^2)$, $(2),(1^4)$, $(2,1^2),(1^6)$ and $(2,1^4)$.
The root bound excludes $(2,2)$, and $(1^8)$ cannot fit in $m\le7$
coordinates. There are also $4\binom m2$ exceptional roots.
Multiplication by $a_r$ and addition of the exceptions proves
\eqref{eq:e8-coordinate-count}. Full column rank, not radial moments
alone, establishes these counts for the fixed cuts.

For the specified $D_5$ correspondence put $H_4=J_4-2I_4$ and
let $P_{5,8}$ exchange coordinates five and eight. The map
\[
 Q=\tfrac12P_{5,8}\operatorname{diag}(H_4,H_4)
\]
permutes the $E_8$ roots; the certificate expresses it as 31 simple
reflections.
It sends $\alpha_1,\ldots,\alpha_5$ to
\[
 -e_1-e_5,\quad e_3+e_4,\quad e_1-e_2,\quad
 e_2-e_3,\quad e_3-e_4,
\]
which span the first five coordinate directions. Thus $Q^{-1}$ sends
the rank-five cut to the specified $D_5$. No ambient Weyl action is
required.
\end{proof}

\begin{table}[ht]
\centering
\caption{Exact fixed coordinate-section counts and rank certificates.
The minors use every orbit column and the specified subset of moment rows.}
\label{tab:e8-coordinate-family}
\begin{tabular}{rrrrr}
\toprule
$k$&Moment rows&Orbit columns&Minor modulo $1000003$&$|X\cap U_k^\perp|$\\\midrule
1&45&25&910777&16815624\\
2&62&34&790744&10663920\\
3&70&39&894205&6688960\\
4&72&41&314211&4149504\\
5&70&39&401115&2553792\\
6&62&34&244188&1593664\\
7&45&25&296071&1110360\\
8&19&13&272057&1092000\\\bottomrule
\end{tabular}
\end{table}

\subsection{What accounts for the improvement}
The comparison in~\cite{crosssections} has $2,545,056$ contacts from
a primitive $\sqrt3D_5$ in $P_{48p}$. Both sections have determinant
972, singleton cluster, $\beta=0$ and permutation-isometric Grams.
The gain therefore lies in the embedding: the specified $E_8$ parent
in $P_{48n}$ forces the larger perpendicular shell, despite identical
intrinsic section data.
